\ssr{ВВЕДЕНИЕ}

Управление реестрами docker-образов – одна из основных технологий в области контейнеризации и DevOps, обеспечивающая надежное хранение, распространение и контроль версий контейнерных артефактов. В условиях повсеместного перехода на микросервисную архитектуру и облачные технологии, вопросы эффективного управления образами становятся критически важными для обеспечения непрерывной интеграции и доставки программного обеспечения.

Актуальность использования систем управления реестрами обусловена необходимостью создания безопасных и отказоустойчивых конвейеров поставки ПО, а также предотвращения уязвимостей в цепочке поставок программного обеспечения.

Одним из наиболее перспективных решений для управления реестрами является система Harbor. Harbor представляет собой открытый реестр контейнерных образов корпоративного уровня, который обеспечивает комплексную безопасность, включая сканирование уязвимостей, цифровую подпись образов и управление доступом на основе ролей. В отличие от базовых реестров, Harbor предоставляет расширенные функции репликации, управления политиками хранения и поддержки различных типов артефактов, что делает его оптимальным решением для крупных организаций \cite{harbor-docs}.

Цель работы -- провести сравнительный анализ современных систем управления реестрами docker-образов на основе платформы Harbor.

Для выполнения поставленной цели необходимо выполнить следующие задачи:
\begin{enumerate}
	\item изучить основные понятия и задачи управления реестрами контейнерных образов;
	\item проанализировать существующие программные решения для управления реестрами;
	\item выделить ключевые критерии для сравнительной оценки систем управления реестрами;
	\item провести сравнительный анализ решений на основе сформулированных критериев;
	\item определить преимущества и ограничения системы Harbor в сравнении с аналогами;
	\item дать оценку поддержки критериев среди сущетсвующих систем.
\end{enumerate}

\clearpage